% !TEX root = ../main.tex
\documentclass[../main.tex]{subfiles}
\begin{document}
\chapter{Conclusion and Outlook} \label{conc}
This thesis focuses on different bottlenecks in geophysical flows and extracts new behavior from fundamental physics. In section \ref{tide-eddy}, we turned to the Boussinesq equation on an $f$-plane and investigated tide-eddy interactions across length scales. We showed that the submesoscale tide-eddy interaction differs significantly from the commonly studied mesoscale ones. At submesoscales, energy interactions are strong enough to form a broadband frequency spectrum signifying strong turbulent mixing. Mesoscale interactions, on the other hand, only excite flows with frequencies of tidal harmonics. We thus uncovered new energy pathways in tide-eddy interactions that can significantly affect ocean model estimates. 

In section \ref{trunc-euler}, we took up the Galerkin-truncated Euler equation where the viscosity in the Navier-Stokes equation was blindly set to zero. An artifact of the finite resolution and truncated wavenumber leads to eventual thermalization of this energy-conserving system. We found that the thermalization timescale is surprisingly resolution-independent. More interestingly, the system shows a constant energy flux across scales without a Kolmogorov $k^{-5/3}$ spectrum. The constant flux appeared in the thermalized modes and was also a direct consequence of a finite truncation wavenumber.  


Section \ref{fn-cstcs} investigated caustics of inertial particles and extended the study from particles infinitely denser than the fluid to all particle densities larger than the fluid. It found similar behavior for caustics across particle-fluid density ratios. High-strain regions trigger caustics, with denser particles requiring a lower strain threshold. It also pointed out the importance of the kind of strain for caustics. Survivor particles, those that encounter regions of very high strain but do not form caustics, originated from regions of extensive strain. They had straighter paths and faster speeds than their caustics counterparts and spent little time in large compressive strain, thus avoiding caustics. Caustics particles had slower speeds and zig-zag trajectories. They remained confined in the compressive strain long enough for caustics. 

Section \ref{act-cstcs} showed another interesting result: caustics of inertialess self-propelled active particles are similar to those of inertial ones. Large strain regions trigger these collisions. The mobility of these particles acts as an effective mass in their dynamics. We pointed out the similarities and differences of their evolution from inertial particles in both point-vortex and $2$D turbulence. 
However, the studies in sections \ref{fn-cstcs} and \ref{act-cstcs} considered collisions between particles of the same size. We went beyond monodisperse particle collisions in section \ref{lck-drp} and developed reduced models to capture the growth of water droplets in atmospheric clouds. The model avoided the cost of checking for pair particle collisions using a hybrid Euler-Lagrange approach. We then investigated how the initial size of the lucky droplets, ones that detach significantly from the background flow and collect other tiny droplets, affects their relative growth rate. We found that smaller lucky droplets grow faster relative to their initial size in $2$D but not in $3$D. And, the mean dissipation rate of turbulent clouds does not modify the gravitationally driven settling. 

This research has given me expertise in asymptotic expansions, eigenfunction decompositions, singular perturbation theory, scaling analysis, and the statistical mechanics of turbulence in geophysical flows. My work also required high-performance computation, for which I developed GPU and MPI-parallelized pseudo-spectral codes for simulating turbulence and the suspended particles. These codes scale well on HPC clusters and are publicly available on GitHub \cite{githubrepo}. 

\end{document}
